\begin{table}[htbp]
\centering
\footnotesize
\caption{Coeficientes con y sin ajustar por la incidencia (misma muestra y mismos pesos)}\label{tab:incidencia}
\begin{tabularx}{\linewidth}{X r r r}
\toprule
Término & Con incidencia & Sin incidencia & Cambio (\%) \\
\midrule
Edad mediana & \num{-0.788} & \num{-1.240} & \num{-57.3} \\
Tamaño medio del hogar & \num{-12.875} & \num{-21.614} & \num{-67.9} \\
Bachillerato (18–24) & \num{0.241} & \num{0.326} & \num{35.2} \\
Grado universitario (25+) & \num{-1.054} & \num{-0.916} & \num{13.1} \\
Desempleo (16+) & \num{0.366} & \num{0.885} & \num{141.7} \\
Solo seguro público & \num{0.892} & \num{0.884} & \num{-0.9} \\
Otra raza & \num{-0.660} & \num{-1.127} & \num{-70.9} \\
Natalidad & \num{-0.381} & \num{-0.857} & \num{-124.6} \\
Región: Noreste & \num{-15.683} & \num{-9.303} & \num{40.7} \\
Región: Medio Oeste & \num{-4.327} & \num{-4.934} & \num{-14.0} \\
Región: Oeste & \num{-14.131} & \num{-17.237} & \num{-22.0} \\
Bachillerato (25+) & \num{0.354} & \num{0.474} & \num{34.1} \\
\bottomrule
\end{tabularx}
\par\smallskip{\footnotesize\raggedright $R^2$ con incidencia: \num{0.558}; sin ella: \num{0.438}. Un cambio grande indica que la incidencia es confusora (o mediadora) de esa asociación.\par}
\end{table}
